\subsection{Supersonic \acrfull{LES}\label{sec:supersonicLES}}
The results presented below were obtained using the turbulent solver in
\gls{SU2} implemented by Molina \cite{Molina2017}. \\ 

The mesh and dimensions for case $L2$ are shown in \cref{fig:supersonicMeshLES}.
The \gls{AO} properties were collected along a beam propagating in the normal
direction ($\hat{y}$), spanning from $-0.85\Units{m}$ to $0.0\Units{m}$. The
beam had a width in the streamwise direction ($\hat{x}$) ranging from
$1\Units{m}$ to $4\Units{m}$.
\begin{figure}[H]
    \makebox[\textwidth][c]{
        %trim = left, bottom, right top
        \includegraphics[width=1\textwidth]{chapter4/supersonicLES/supersonicMesh}}
    \caption{Mesh for $L2$}\label{fig:supersonicMeshLES}
\end{figure}

The Mach contour plot is shown in \cref{fig:supersonicMachContourLES}.
The region where the \gls{AO} properties would be analyzed lie in the supersonic regime, which is consistent with expectations, since across a shock wave the post-shock Mach number is always lower than the pre-shock Mach number.
\begin{figure}[H]
    \makebox[\textwidth][c]{
        %trim = left, bottom, right top
        \includegraphics[width=1\textwidth,trim=0.8cm 15cm 0.8cm
        10cm,clip]{chapter4/supersonicLES/Mach_07100}}
    \caption{Mach contour for $L2$}\label{fig:supersonicMachContourLES}
\end{figure}

The velocity contour plot is shown in \cref{fig:supersonicVelocityContourLES}.
\begin{figure}[H]
    \makebox[\textwidth][c]{
        %trim = left, bottom, right top
        \includegraphics[width=1\textwidth,trim=0.8cm 15cm 0.8cm 10cm,clip]{chapter4/supersonicLES/Velocity_07100}}
    \caption{Velocity contour for $L2$}\label{fig:supersonicVelocityContourLES}
\end{figure}

The temperature contour plot is shown in \cref{fig:supersonicTemperatureContourLES}.
\begin{figure}[H]
    \makebox[\textwidth][c]{
        %trim = left, bottom, right top
        \includegraphics[width=1\textwidth,trim=0.8cm 15cm 0.8cm
        10cm,clip]{chapter4/supersonicLES/Temperature_07100}}
    \caption{Temperature contour for
    $L2$}\label{fig:supersonicTemperatureContourLES}
\end{figure}

\iffalse
The pressure contour plot is shown in \cref{fig:supersonicPressureContourLES}.
\begin{figure}[H]
    \makebox[\textwidth][c]{
        %trim = left, bottom, right top
        \includegraphics[width=1\textwidth,trim=0.8cm 15cm 0.8cm
        10cm,clip]{chapter4/supersonicLES/Pressure_07100}}
    \caption{Pressure contour for $L2$}\label{fig:supersonicPressureContourLES}
\end{figure}
\fi

The density contour plot is shown in \cref{fig:supersonicDensityContourLES}.
\begin{figure}[H]
    \makebox[\textwidth][c]{
        %trim = left, bottom, right top
        \includegraphics[width=1\textwidth,trim=0.8cm 15cm 0.8cm
        10cm,clip]{chapter4/supersonicLES/Density_07100}}
    \caption{Density contour for $L2$}\label{fig:supersonicDensityContourLES}
\end{figure}

The polarizability contour plot is shown in \cref{fig:supersonicPolarizabilityContourLES}. It was calculated using the Kerl method \cite{Kerl1992}, implemented in the \gls{HAOT} package. The method is described in \cref{eq:kerlFitting}. As expected, Kerl's extrapolation depends on temperature; consequently, the polarizability closely follows the temperature field.
\begin{figure}[H]
    \makebox[\textwidth][c]{
        %trim = left, bottom, right top
        \includegraphics[width=1\textwidth,trim=0.8cm 15cm 0.8cm
        10cm,clip]{chapter4/supersonicLES/kerl_polarizability_07100}}
    \caption{Kerl-Polarizability contour for $L2$}\label{fig:supersonicPolarizabilityContourLES}
\end{figure}

The dilute-medium formulation of the index of refraction is shown in
\cref{fig:supersonicDiluteIndexContourLES}. It was calculated using
\cref{eq:diluteRefractiveIndex}, which calculates the index of refraction as a function of density and temperature.
\begin{figure}[H]
    \makebox[\textwidth][c]{
        %trim = left, bottom, right top
        \includegraphics[width=1\textwidth,trim=0.6cm 15cm 0.8cm
        10cm,clip]{chapter4/supersonicLES/dilute_index_of_refraction_07100}}
    \caption{Dilute-medium formulation of the index of refraction contour for $L2$}\label{fig:supersonicDiluteIndexContourLES}
\end{figure}

The dense-medium formulation of the index of refraction is shown in
\cref{fig:supersonicDenseIndexContourLES}. It was calculated using
\cref{eq:denseRefractiveIndex}, which calculates the index of refraction as a
function of density and temperature. It does not seems to be a noticeable
difference between the dilute-medium formulation and the dense-medium
formulation. It seems that the index of refraction closely follows the density
field.
\begin{figure}[H]
    \makebox[\textwidth][c]{
        %trim = left, bottom, right top
        \includegraphics[width=1\textwidth,trim=0.6cm 15cm 0.8cm
        10cm,clip]{chapter4/supersonicLES/dense_index_of_refraction_07100}}
    \caption{Dense-medium formulation of the index of refraction contour for
    $L2$}\label{fig:supersonicDenseIndexContourLES}
\end{figure}

The \gls{OPL} (computed using the dilute-medium formulation) for a given time step is shown in \Cref{fig:supersonicOPL633}, and it was obtained from \cref{eq:OPL}, as implemented in the \gls{HAOT} package. A sharp increase in the \gls{OPL} occurs at approximately $-0.4\Units{m}$, which corresponds to the region just outside the shock. In this region, the density is significantly higher. Because the index of refraction is related to density through the \gls{GD} relation, an increase in density directly increases the local index of refraction. This, in turn, results in a larger optical path length and consequently produces a stronger contribution to the \gls{AO}.
\begin{figure}[H]
    \makebox[\textwidth][c]{
        %trim = left, bottom, right top
        \includegraphics[width=0.6\textwidth]{chapter4/supersonicLES/opl_07100}}
    \caption{\gls{OPL}, dilute-medium formulation for a wavelength of
    $633\Units{nm}$}\label{fig:supersonicOPL633}
\end{figure}

The \gls{OPD} (computed using the dilute-medium formulation) for a given time
step is shown in \Cref{fig:supersonicOPD633}, and it was obtained from
\cref{eq:OPD}, as implemented in the \gls{HAOT} package. From
\cref{fig:supersonicOPL633} can be observed that the spatial standard deviation
is much smaller than the spatial average \gls{OPL}. When this is the case, the
wavefront is bassically flat and the \gls{OPD} does not revel additional information.
\begin{figure}[H]
    \makebox[\textwidth][c]{
        %trim = left, bottom, right top
        \includegraphics[width=0.6\textwidth]{chapter4/supersonicLES/opd_07100}}
    \caption{\gls{OPD}, dilute-medium formulation for a wavelength of
    $633\Units{nm}$}\label{fig:supersonicOPD633}
\end{figure}
